\subsection{Current evaluation parents sharpen the matched-design target}

The 2026 evaluation landscape makes a broad ``agent quality'' comparison inadequate. SciConBench evaluates 9.11K scientific-conclusion questions under a clean-room harness and reports that leakage-sensitive unrestricted settings can substantially overestimate factual synthesis; the best clean-room system in that study remains low on factual precision/recall composition \cite{jung2026scicon}. SciAgentArena instead supplies approximately two hundred interactive real-research tasks with stepwise verification and finds a capability frontier: agents are comparatively effective in well-specified analysis workflows while remaining weaker on open-ended exploration and novel insight generation \cite{liu2026sciagentarena}. These results motivate sealed evidence universes, task-level verification and regime-stratified reporting rather than a single average score.

Evidence binding is another inherited parent mechanism. GAVEL requires atomic subclaims to be attached to explicit evidence units and adds deterministic scrutiny of cited identifiers and spans; its ablations support the value of both evidence binding and mechanized validation \cite{xu2026gavel}. RIGOURATE separately operationalizes scientific overstatement by retrieving the evidence associated with a claim and estimating whether the prose exceeds what that evidence supports \cite{james2026rigourate}. RAKL should assimilate these strengths rather than claim novelty for evidence contracts or evidence-proportional claims. The additional Paper-2 target is whether those local contracts remain coherent when the scientific state also has context, relation type, evidence-lineage dependence, mechanism/identification separation, negative history and protected authority transitions.

Finally, a completed research trajectory is not automatically a body of evidence. Zhuang et al. explicitly study trajectory-to-evidence conversion: consequential artifacts are verified before release, executed rounds are qualified after execution, applicability boundaries and provenance are preserved, and later rounds can underperform earlier ones \cite{zhuang2026trajectory}. This is especially relevant to Challenge Learning and Self-RAKL. A method generation, tool trajectory or apparently successful repair earns only proposal/development status until the result is qualified against the frozen target and fresh assurance. Therefore the matched evaluation records the best validated generation and its lineage rather than assuming that the final iteration is epistemically strongest.

Together these parent methods imply a stricter causal design. Paper 2 will separately manipulate evidence access and control architecture, protect task/evaluator identity, score process failures before endpoint averaging, and report capability frontiers by evidence topology. A simpler parent winning in a low-conflict or single-signal regime is not an embarrassment; it is evidence about where the extra epistemic machinery has not earned its cost.
